\chapter{Where Is Humanity Heading?}
\section{Pick Up a Fire-Cracked Bone}
Pick up not a phone or coin, but a palm-sized cattle shoulder blade, yellowed and rounded at the edges. Museums' Shang-dynasty cases contain many such bones.

Turn it over. One side has rows of round or oval hollows deliberately cut with tools. The other bears fine cracks like a dried pond's floor, with lines of ancient writing between them: oracle-bone inscriptions, roughly three thousand years old.

What was it for? Asking the future. Faced with uncertain battles, royal childbirth, or harvests, people heated bones until they cracked, read answers in the patterns, then inscribed questions and answers.

Do not laugh too quickly. We consult weather forecasts, predicted university-admission thresholds, economic outlooks, and lists of occupations AI might replace within ten years. The method changes; the impulse does not. From cracked bone to predictive notification runs one theme: \textbf{humans alone among animals lose sleep over their own tomorrow}. Others store food and avoid predators; humans ask where they will be in ten years, whether the world will remain in a hundred, and where humanity is heading.

This chapter asks the largest question, traveling from bone through temples, fields, laboratories, and networks before returning it to you.

\section{Asking the Future Three Thousand Years Ago}
Enter a scene with me.

Time: an early morning around the thirteenth century BCE. Place: Yin, Shang's capital near present-day Anyang, Henan. Dawn carries smells of wood smoke and livestock. Several people gather around an unlit fire in a palace courtyard.

An important question awaits: Fu Hao, King Wu Ding's wife, is near childbirth. Other inscriptions reveal that this unusual queen commanded armies numbering over ten thousand. Today, however, they ask not about war but whether delivery will go well.

The ceremony's specialist, a diviner responsible for questioning the future, is named Que, a name inscribed so we can still read it three millennia later. He takes a prepared cattle scapula, recites the question, and presses a glowing branch into a hollow.

After seconds, a faint crack sounds.

Fissures spread like tiny frozen lightning. Onlookers lean closer. Que records the question; the highest-ranking person, King Wu Ding, reads the cracks and pronounces his judgment. Question and judgment are inscribed. Remarkably, a month later, someone returns to add what actually happened.

Notice that this ritual includes question, answer, and \textbf{review}. Verification joins the inscription. It is not merely an empty supernatural performance, but an institution for managing uncertainty: ask seriously, record, check afterward. Its structure startlingly resembles modern forecast publication and accuracy review.

Fu Hao's delivery outcome survives on the bone. We will read it closely later; it reveals more than you might expect.

\section{Can the Future Be Known?}
State the problem before answering.

The bone confronts us with a contradiction. Predictive ability has advanced enormously. Shang kings burned bones for harvests; we forecast typhoon tracks a week ahead. Comets once suggested dynastic doom; we calculate Halley's next return to the year. Eclipses, tides, population, and vaccine protection can be predicted fairly accurately.

Yet predictive failure is equally spectacular, often concerning what matters most: war, epidemics, transformative technology, falling regimes. Afterward, everyone sees signs; beforehand, almost nobody predicts correctly. Even leading experts perform much worse on such events than they imagine, as we will see.

The question becomes a chain:

\textbf{Is the future already written but unread, or is there no script and we make it step by step?}

If fixed, every effort, choice, and dispute merely recites lines, differing only in whether we know the next one. Many secretly believe this, from incense offerings and horoscopes to destiny and inevitable historical laws: a script already exists.

If open, the burden grows. Every present choice genuinely changes the world; nobody can own your future responsibility. Leaving later matters for later fails because later consists of countless nows.

Before proceeding, establish a signpost: \textbf{history has no prewritten destination}. This does not deny regularities: water flows downward without a script. Nor does it prohibit discussing trends, probabilities, and scenarios. It denies one guaranteed endpoint, whether universal harmony, history's end, or inevitable extinction. The chapter's evidence will support this.

Choose provisionally: if a machine could accurately describe your life thirty years hence, would you want to know?

\section[Kings, Farmers, Oracles, and the Future]{Kings, Farmers, and Oracles: Betting on the Future}
Hear three voices: rulers, ordinary people, and prediction's professionals.

\textbf{First: a king wants to purchase certainty.}

Travel west from Yin and forward to the sixth century BCE. Lydia's King Croesus, in present-day Turkey, was wealthy enough to lend his name to the phrase rich as Croesus. Unsure about attacking Persia, he tested oracles, according to Herodotus's Histories, Book I. Envoys asked different temples on the same day what the king was doing. Delphi's priestess correctly described turtle and lamb cooking in a bronze pot, his deliberately strange test. Trusting Delphi, he made lavish offerings and asked whether to attack Persia.

The oracle, as Herodotus recounts, broadly said that crossing the Halys to attack would destroy a great empire.

Delighted, Croesus marched and suffered catastrophe. The destroyed empire was his own.

Beyond a joke about superstition, notice due diligence and expensive expert consultation. The problem was wording that fit either outcome: victory destroys Persia; defeat destroys Lydia. \textbf{Ambiguity keeps prophecy alive}. Today's online masters use the same structure, blocking no outcome and claiming success afterward. Your first defense: a prediction correct under every outcome says nothing.

\textbf{Second: ordinary people stake entire lives on a forecast.}

Jump to rural New England in the 1840s. Baptist preacher William Miller spent years studying scripture and calculated Christ's return and the world's end between 1843 and 1844. Tens of thousands or more believed him. Some sold farms, left jobs, or abandoned harvests deemed useless. October 22, 1844 became the deadline. Many waited overnight on hillsides and rooftops in white clothing.

Dawn came. Nothing happened.

This became the Great Disappointment. Do not turn believers into jokes. Farmers, shoemakers, and teachers sincerely staked everything on an assertion about the future. Sold fields and lost crops were real. Some broke down or became poor; others formed new churches lasting today. Ask not why they were foolish but \textbf{why a sincere prophecy believed by thousands could still be wrong}. The answer is hard: conviction's strength is not evidence's strength. Certainty and reliability differ, for apocalyptic predictions as for ever-rising house prices or AI destroying humanity within three years.

\textbf{Third: the recorder's honesty.}

Return to Que and his colleagues. They were not decision-makers; kings interpreted cracks. Diviners recorded questions, judgments, and subsequent fulfillment. An unexpectedly unsuperstitious seed lay here: \textbf{archives}. Records let later readers distinguish successes from failures. Verification inscriptions turned ritual into a by-product, history itself. Receipts for questioning the future became our most precious Shang sources. Humanity strained to know tomorrow and left its greatest value in honest records of yesterday.

Kings purchase certainty, farmers stake everything, recorders preserve questions and answers. Remember these three postures; we repeatedly choose among them before the future.

\section[Trends, Scenarios, and Surprises]{Thinking Bridge: Trends, Scenarios, and Surprises}
Can we analyze the future with a portable tool rather than feeling? In the later twentieth century, futures studies developed with military, corporate, and government funding. Without crystal balls, it offered a useful distinction: \textbf{sort future claims into three kinds}.

\textbf{First: trends, paths already followed and likely to continue awhile.}

Demography, literacy, life expectancy, energy prices, and urbanization are slow variables with riverlike momentum. Use them by asking where curves may bend, not simply extrapolating today's increment into tomorrow. In 1798, Thomas Malthus wrote in the first chapter of An Essay on the Principle of Population: ``Population, when unchecked, increases in a geometrical ratio. Subsistence increases only in an arithmetical ratio.'' Population doubles 1, 2, 4, 8 while food climbs 1, 2, 3, 4, making hunger seem inevitable. Such trends recurred for two centuries, including 1960s bestsellers predicting unavoidable famine. Yet fertilizer, improved crops, and machinery let food outgrow population. Malthus's arithmetic was not the mistake; treating contemporary agricultural technology as fixed was. \textbf{Extrapolation's greatest enemy is the trend's ability to turn}.

A widely repeated caution, possibly embellished in detail, says late-nineteenth-century London and New York planners feared horse numbers would bury streets in manure. The crisis vanished not through better cleaning but automobiles, which invalidated the curve and its axes. Do not dismiss trends; add unless the track changes whenever someone says at this rate.

\textbf{Second: scenarios, preparing several futures instead of guessing one.}

A twentieth-century practice improved on prediction by preparing multiple coherent scripts. In the early 1970s, Shell planners including Pierre Wack avoided betting everything on stable oil prices. They explored coordinated producer price rises and prepared responses. When the 1973 crisis multiplied prices, Shell responded more calmly than many rivals. Scenario planning promises not which script becomes real, but preparedness whichever does. It honestly admits uncertainty without accepting unreadiness.

\textbf{Third: wildcards, the unscheduled card.}

Some events escape trends and scenarios. Black swans are rare shocks scarcely anticipated but afterward deemed obvious. Gray rhinos are huge repeatedly warned risks people keep discussing as they charge. The early-twenty-first-century global pandemic exemplifies a gray rhino: epidemiologists warned for decades of another influenza-scale outbreak while reports sat in government cabinets. The tool does not teach exact surprise prediction, which would remove surprise, but reserves: food, money, redundancy, alternatives. Systems staking everything on optimal forecasts are least resilient to the unexpected.

Sort claims about guaranteed careers, rising property, or technology ending everything into these drawers. Trend, scenario, or concern about surprise? What evidence, and what conditions would falsify it? Most absolute prophecies cannot even enter the first drawer.

\section{Reading an Ancient Question and Answer}
Return to Fu Hao's oracle, among the most famous surviving Yin inscriptions. Here is the broad scholarly reading, with individual character identifications still discussed:

\begin{quote}
Divined on jiashen; Que asked: Fu Hao's delivery---auspicious? The king read: if on a ding day, auspicious; if on a geng day, greatly favorable. Thirty-one days later, on jiayin, she delivered. Not auspicious: a girl.
\end{quote}
In today's language:

On jiashen, Que asked whether Fu Hao would deliver safely. Wu Ding read the cracks: a ding-day birth would go well, a geng-day birth better still. Thirty-one days later, on jiayin, she gave birth: not favorable, a girl.

These thirty-odd characters condense almost every chapter theme and deserve word-by-word attention.

First, the question concerns a real, risky birth. Without modern obstetrics, childbirth threatened queens too. The anxiety equals any expectant parents' today, palpable across three millennia.

Second, the answer marks ding and geng as favorable windows. Predictions have elasticity: expected births occupy a range, and two dates in it are labeled lucky. Outcomes readily appear fulfilled. This does not prove deliberate royal fraud, but ambiguity built into even solemn ancient prophecy.

Most important is the added outcome: \textbf{it did not work}. Neither predicted date matched, and the girl was labeled unfavorable according to contemporary values. That painful ranking records their thinking, not what ours should be. They could have omitted failure or preserved only successes, yet inscribed the incorrect prediction in royal archives.

Why does that matter? \textbf{Prediction can improve only when paired with review}. Philip Tetlock's famous long-running study scored hundreds of experts' forecasts. Broadly, without scoring and checking, even elite long-term political predictions do scarcely better than coin flips; careful records enable measurement, comparison, and improvement. The ancient recorder who carved not auspicious left a testable answer sheet. We know the oracle failed because its users honestly recorded failure.

\section{The Twentieth Century: Better, and Why?}
Apply the tools to a real question. Over the two hundred years since the twentieth century began, global life expectancy rose from thirty-odd years to around seventy; literacy became common rather than privileged; extreme poverty fell from roughly ninety to ten percent. WHO declared smallpox eradicated in 1980. These magnitudes rest on repeated historical and multinational statistical checks, not merely opinions.

Three explanations compete. Separate what happened, the numbers, from why, the forthcoming explanations; participants' understandings come later, and evaluation remains yours.

\textbf{Explanation one: improved institutions and ideas.} Smallpox eradication required humanity's first globally coordinated public-health campaign: vaccines, tracing, isolation, international organization. Compulsory schooling, women's education, and widespread antibiotics also rest on institutions. Yet institutions and ideas can reverse. Some countries proudest of science legalized eugenics and forcibly sterilized supposedly unfit citizens. In 1927, the US Supreme Court upheld Carrie Buck's sterilization in Buck v. Bell. Justice Holmes's words, ``Three generations of imbeciles are enough,'' still chill. Buck was later shown not to fit that label, but to be a poor ordinary girl crushed by institutions. A progressive age can commit barbarity in science's name. The account's remaining gap is what sustains institutions against reversal.

\textbf{Explanation two: energy and technology's dividend.} Credit coal, oil, electricity, and fertilizer. In the early twentieth century, Fritz Haber developed ammonia synthesis from air, enabling mass fertilizer production underlying billions of meals. Yet his other life complicates the evidence: he oversaw wartime chlorine deployment, personally supervising the first large-scale poison-gas attack at Ypres in 1915. His wife, also a chemist, died by suicide soon afterward. One hand fed billions; another opened chemical warfare's Pandora's box. \textbf{Technological progress does not automatically bring moral progress}. Technology amplifies users' intentions: ammonia fertilizes or supplies explosives, nuclear energy powers cities or destroys them, the internet distributes knowledge or mass rumors. Expecting technology's own morality resembles expecting a knife to choose vegetables over victims.

\textbf{Explanation three: much was luck.} The 1962 Cuban Missile Crisis brought thirteen days near nuclear war. Later declassified material describes a Soviet submarine under American depth-charge pressure on October 27, without Moscow contact. Officers disputed firing a nuclear torpedo: required unanimity among three was blocked by Vasili Arkhipov against two approvals. No torpedo launched. The later twentieth century's absence of nuclear war therefore contains substantial luck. Survivors mistake disaster's nonoccurrence for impossibility, a survivorship bias. This honest account discounts optimism, but if everything is luck, what value has effort? The crisis was controlled precisely because someone resisted procedure. Luck and effort intertwine.

Before criticizing, \textbf{make optimism's strongest case}. Selling pessimism is also bias and often more marketable: gloom looks profound, optimism naive. Improved longevity, literacy, nutrition, and infant survival changed billions of real lives, not merely statistics. Progress deniers often belong to comfortable generations untouched by hunger or smallpox. Pessimism also costs action: young people expecting inevitable decline may make it self-fulfilling. The proper stance is therefore \textbf{both pessimism and optimism require evidence}. Confidence deciding victory belongs to debate contests, not understanding the world.

\section[Does History Have a Direction?]{Deeper Challenge: Does History Have a Direction?}
Our deepest question moves beyond what happens next to \textbf{whether history itself has direction}. Is society heading somewhere? Across two millennia, three broad answers have distinguished advocates and cross-cultural relatives.

\textbf{Position one: history cycles.} Seasons, lives, and dynasties make cycles intuitive worldwide. Chinese readers know Romance of the Three Kingdoms' opening:

\begin{quote}
The world's great tendency is this: long divided, it must unite; long united, it must divide.
\end{quote}
Division and unity recur without endpoint. Fourteenth-century North African thinker Ibn Khaldun refined the idea in the Muqaddimah: cohesive groups establish rule, comfortable descendants decay, and after roughly three generations fresh cohesive groups replace them. Cyclical history is humble, having seen many claims that this time differs end similarly. Its weakness is irreversible change. Smallpox eradication does not cycle like kingdoms, nor literacy return to oracle-bone levels. Some doors cannot close after passage.

\textbf{Position two: history advances straight toward progress.} Condorcet's life makes this position tragicomic. Hunted during the French Revolution's bloodiest period, he hid in a Paris room and wrote Sketch for a Historical Picture of the Progress of the Human Mind, arguing reason and freedom would advance without limit and abolish ignorance and tyranny. He finished, fled, was arrested, and died in prison in 1794 under the tyranny progress supposedly would eliminate. The book appeared posthumously. Cynics say progress could not save its author; admirers say belief beneath the guillotine's shadow refused to be bought by present darkness. Hold both readings. Progress has real data behind it, but can quietly turn recent improvements into everything automatically improving, then automatic into inevitable. Consider Norman Angell's 1910 The Great Illusion. Its strong economic argument held that interdependent finance and trade made great-power war ruinous and unprofitable. Many read it as proving major war impossible. Four years later, trenches consumed a generation. The lesson is not reasoning's uselessness, but its inability to fully calculate fear, vanity, misjudgment, and chance, history's regular guests.

\textbf{Position three: structures without a script.} Complex-systems thinking offers a later, harder account. Weather follows physical laws, yet precise conditions two weeks ahead escape forecasts because interacting factors amplify tiny initial differences. The butterfly effect means not that butterflies straightforwardly create storms, but that \textbf{small differences in initial conditions grow exponentially}. Human societies are more complex: billions of decisions, technology, pathogens, weather, and one person's thought, such as Arkhipov's no, alter direction. \textbf{History can be understood but not predicted, with patterns but no script}. Looking backward, steps make sense; looking forward, forks exceed our maps.

These are not lottery choices. Cycles capture repetition but not irreversible shifts. Progress captures improvements but not reversals and luck. Complexity captures unpredictability but must not become a couch for inaction. Complete the opening signpost: \textbf{no script does not mean no responsibility. Without a script, actors genuinely answer for their lines}.

\section{Your Future Is for Sale in Your Pocket}
The ancient question performs daily on your screen.

Platforms flood you with professions doomed in ten years, AI replacing office workers, halving house prices, and civilization-changing technologies. Sorted into our drawers, most prove neither serious trends, scenarios, nor honest warnings, but \textbf{emotional commodities}. Algorithms learn that fear and hope attract clicks, manufacturing limitless doom and sudden riches for tomorrow's anxious people. Croesus paid sacrifices for one oracle; you receive a hundred freely. Yet what is free is not prophecy but you: your anxiety is somebody else's priced product.

Examine two contemporary disputes.

First, AI. Capabilities recently advanced beyond many experts' expectations. In 2023, hundreds of researchers and business leaders signed a one-sentence statement placing mitigation of AI extinction risk alongside pandemics and nuclear war as a global priority, including respected scientists. Equally experienced experts dispute the severity, emphasizing immediate employment, bias, and misinformation harms. Separate facts from interpretations and open judgments. They dispute extrapolation for something unprecedented, with unknown assumptions about continuing capability growth and timely institutions. This is an open question, not a settled issue merely awaiting allegiance.

Second, climate. Pessimists have genuine warming, retreating-glacier, and extreme-weather evidence. Optimists have genuine renewable cost reductions exceeding an order of magnitude over little more than a decade, electric cars becoming ordinary, and successful atmospheric cooperation through the 1987 Montreal Protocol's phased elimination of ozone-depleting chemicals. The ozone hole entered long recovery. This success receives too little attention because it lacks frightening clicks. Yet it offers vital evidence: \textbf{the future is negotiated, legislated, built, and sustained, not merely predicted}.

We can now fulfill the opening promise: four things matter equally under uncertainty.

\begin{itemize}
\item \textbf{Institutions}: turn goodwill and foresight into executable arrangements, from Montreal and vaccination systems to flood plans and school fire drills. They are the future that keeps working when we forget.
\item \textbf{Memory}: place lessons in public archives, from Que's unfavorable result to Buck's case and pandemic reviews. Without this foundation, every generation pays tuition from scratch.
\item \textbf{Imagination}: rehearse possibilities beyond reality, through More's Utopia, Wells's The Time Machine, science fiction, and planning models. Living futures in stories helps choose and avoid them. Peoples unable to imagine receive others' imagined futures.
\item \textbf{Public discussion}: let futures collide in votes and debate rather than private decisions by a few. Classroom councils, community hearings, and public online arguments shape which future society receives through their quality.
\end{itemize}
None predicts the future. Together they determine whether a charging gray rhino meets a prepared community or isolated people scrolling phones.

\section[Spending Today for Future Strangers]{For You: Spending Today for People You Will Never Meet}
Return to the bone.

Three thousand years ago, a dynasty solemnly recorded Fu Hao's delivery question and failed prediction. It could not know that bone would reveal its deepest anxiety and honest recording to us. Our turn comes now. What will people three thousand, or merely three hundred, years hence read in our oracle bones: seabed plastics, exhausted and regrown forests, compromised treaties, ineradicable data-center disputes, or something else?

Beneath this lies a question already awaiting you:

\textbf{Is paying today's costs for a future you probably will not see worthwhile?}

Climate action charges now while chiefly benefiting unborn people decades later. Wetland protection, libraries, basic research, and better rules for unknown successors share that structure.

Weigh both sides before answering.

One says future people cannot vote for or thank you and may never exist. Discounted invisible benefits approach zero. Every generation has difficulties; why sacrifice for descendants while today's illness and rent remain? Spending on the visible present is honesty, not selfishness.

The other says every relatively clean breath, vaccine, inherited book, and sound bridge was paid for by strangers who never met you. You inherit an enormous future estate. Civilization is a relay through time: predecessors plant, descendants shade themselves, and nobody sees the whole course. If every generation spent only on itself, the baton would fall at the first exchange, leaving no continuing generations.

Does an unborn person have a claim on today? What do you owe the future? Or is debt the wrong word, and freedom the right one: freedom to become an ancestor worth remembering in three thousand years?

There is no standard answer. But occasionally imagine a young person like you, millennia later, standing before a case and trying to read you through something your age left behind.
